\section*{अध्याय \ref{ch:g4:conductors-insulators} --- चालक और विद्युतरोधी}

\begin{solution}{exo:g4:conductors-insulators:1}
\omterm{def:g4:conductors-insulators:conductor}{चालक} बिजली को अपने में से गुज़रने देता है; \omterm{def:g4:conductors-insulators:insulator}{विद्युतरोधी} उसे रोक देता
है। \omterm{def:g3:first-electric-circuit:bulb}{बल्ब} \omterm{def:g4:conductors-insulators:conductor}{चालक} जलाता है --- वही फंदा बंद करता है।
\end{solution}

\begin{solution}{exo:g4:conductors-insulators:2}
\omterm{def:g4:conductors-insulators:conductor}{चालक}: इस्पात का चम्मच, रसोई की पन्नी, ताँबे का सिक्का।
\omterm{def:g4:conductors-insulators:insulator}{विद्युतरोधी}: कॉर्क, गिलास, प्लास्टिक का ब्लॉक।
\end{solution}

\begin{solution}{exo:g4:conductors-insulators:3}
दोनों खुले सिरे सीधे आपस में छुआओ: \omterm{def:g3:first-electric-circuit:bulb}{बल्ब} जलना ही चाहिए। अगर ख़ुद
जाँच-यंत्र ख़राब हो, तो हर वस्तु \omterm{def:g4:conductors-insulators:insulator}{विद्युतरोधी} लगेगी --- यह जाँच बेगुनाह
पर इल्ज़ाम लगने से बचाती है।
\end{solution}

\begin{solution}{exo:g4:conductors-insulators:4}
या तो वस्तु \omterm{def:g4:conductors-insulators:insulator}{विद्युतरोधी} है --- या जाँच लापरवाह थी: तार का सिरा वस्तु
के नंगे पदार्थ पर कसकर दबने के बजाय रंग, वार्निश या मैल को छू रहा हो
(या कोई जोड़ ढीला पड़ गया हो)।
\end{solution}

\begin{solution}{exo:g4:conductors-insulators:5}
ताँबा चालन करता है: वही लैंप तक धारा की सड़क है। प्लास्टिक रोक लगाता
है: वह धारा को तार के भीतर और उँगलियों को सुरक्षित बाहर रखता है।
\end{solution}

\begin{solution}{exo:g4:conductors-insulators:6}
धातु का नन्हा पुल \omterm{def:g4:conductors-insulators:conductor}{चालक} है --- \omterm{ex:g3:first-electric-circuit:switch}{स्विच} चालू होने पर उसे ही फंदा बंद
करना है। खोल \omterm{def:g4:conductors-insulators:insulator}{विद्युतरोधी} है --- उँगलियाँ दिन भर उसी को दबाती हैं।
पुल के बिना \omterm{def:g1:light-and-shadows:source}{प्रकाश} नहीं; खोल के बिना सुरक्षा नहीं।
\end{solution}

\begin{solution}{exo:g4:conductors-insulators:7}
हाँ --- एल्यूमिनियम धातु है, और हर धातु \omterm{def:g4:conductors-insulators:conductor}{चालक} है। चेतावनी यह है:
चुंबकीय होना और \omterm{def:g4:conductors-insulators:conductor}{चालक} होना अलग-अलग प्रतिभाएँ हैं; \omterm{def:g1:magnets:magnet}{चुंबक} का
सिर्फ़-लोहा वाला नियम चालन के बारे में कुछ नहीं कहता।
\end{solution}

\begin{solution}{exo:g4:conductors-insulators:8}
लकड़ी (और रंग) \omterm{def:g4:conductors-insulators:insulator}{विद्युतरोधी} है, पर भीतर की ग्रेफ़ाइट सलाख़ \omterm{def:g4:conductors-insulators:conductor}{चालक} है ---
एक दुर्लभ अधातु \omterm{def:g4:conductors-insulators:conductor}{चालक}, जो नोक से नोक तक पेंसिल के भीतर छिपा चलता है।
\end{solution}

\begin{solution}{exo:g4:conductors-insulators:9}
\omterm{def:g4:conductors-insulators:insulator}{विद्युतरोधी} इंसान की रक्षा कर रहा है: अगर औज़ार की धातु कभी बिजली
वाली किसी चीज़ को छू ले, तो मोटी प्लास्टिक मूठ धारा को हाथ तक आने से
रोक देती है।
\end{solution}

\begin{solution}{exo:g4:conductors-insulators:10}
धारा को पूरी सड़क पार करनी पड़ती है। ज़ंजीर में धातु की एक कड़ी धातु
की दूसरी कड़ी को छूती है: यानी पूरी धातु की सड़क। मनकों वाली माला में
एक धातु के मनके और अगले मनके के बीच सूती गाँठ का टुकड़ा पड़ता है ---
जो \omterm{def:g4:conductors-insulators:insulator}{विद्युतरोधी} है --- इसलिए सड़क हर मनके पर टूटी हुई है: इतने \omterm{def:g4:conductors-insulators:conductor}{चालक}
होने पर भी अँधेरा।
\end{solution}

\begin{solution}{exo:g4:conductors-insulators:11}
नम त्वचा सूखी त्वचा से कहीं बेहतर चालन करती है, इसलिए गीले हाथ धारा
को शरीर के भीतर से एक आसान सड़क दे देते हैं। \omterm{def:g3:first-electric-circuit:battery}{बैटरी} वाला जाँच-यंत्र
इसलिए सुरक्षित है कि \omterm{def:g3:first-electric-circuit:battery}{बैटरी} का धक्का बहुत हल्का है; दीवार के सॉकेट का
धक्का सैकड़ों गुना ताक़तवर है --- नम त्वचा के रास्ते वह जानलेवा हो
सकता है, इसीलिए नियम \omterm{def:g3:first-electric-circuit:battery}{बैटरी} के खेल की छूट देते हैं और सॉकेट को हमेशा
मना करते हैं।
\end{solution}
